% =====================================================================
%  Задание на ВКР (структура по образцу референсной реализации
%  пользователя; содержание адаптировано к настоящей теме).
%  Страницы: 1) шапка + УТВЕРЖДАЮ + Направление + ЗАДАНИЕ + Тема +
%  исполнители + г. Москва 2026;
%  2+) longtable с разделами Цель/Задачи/...+/Календарный план +
%  даты выдачи/срока + подписи.
% =====================================================================
\thispagestyle{plain}
\fontsize{12pt}{14pt}\selectfont

% ===== Шапка =====
\begin{center}
Министерство высшего образования и науки Российской Федерации\\[2pt]
НАЦИОНАЛЬНЫЙ ИССЛЕДОВАТЕЛЬСКИЙ ЯДЕРНЫЙ УНИВЕРСИТЕТ <<МИФИ>>
\end{center}

\begin{center}
ИНСТИТУТ ИНТЕЛЛЕКТУАЛЬНЫХ КИБЕРНЕТИЧЕСКИХ СИСТЕМ
\end{center}

\vspace{30pt}

% ===== УТВЕРЖДАЮ (справа) =====
{\fontsize{12pt}{14pt}\selectfont
\begin{flushright}
\begin{tabular}{@{}l@{}}
\hfill <<УТВЕРЖДАЮ>>\\
\hfill Зав. каф. №42\\[8pt]
\hfill ~\makebox[2.5cm]{\hrulefill}\,~Епишкина А.В.\\
\hfill <<\makebox[0.7cm]{\hrulefill}>>\,\makebox[3cm]{\hrulefill}\,2026 г.
\end{tabular}
\end{flushright}
}

\vfill
\vspace{18pt}

\begin{center}
\textbf{Направление подготовки} \\
\textbf{10.03.01 <<Информационная безопасность>>}

\vspace{30pt}

\textbf{ЗАДАНИЕ НА ВЫПУСКНУЮ КВАЛИФИКАЦИОННУЮ РАБОТУ}

\vspace{8pt}

\textbf{Тема:}\\
\textbf{Методика обнаружения несанкционированных действий и аномальной \\
активности в корпоративной сети с использованием нейронных сетей}

\vspace{30pt}
\end{center}

\vfill

\noindent
\begin{tabular}{@{}p{4.5cm}@{\hspace{0.4cm}}p{6cm}@{\hspace{0.4cm}}l@{}}
Студент группы Б22-505 & \hrulefill & Титов Д.И.     \\[18pt]
Научный руководитель   & \hrulefill & Когос К.Г.     \\[18pt]
Консультант            & \hrulefill & Ланской Д.П.   \\[18pt]
Рецензент              & \hrulefill & Снегирев А.И.  \\
\end{tabular}

\vfill

% ===== Подвал =====
\begin{center}
г. Москва \\[4pt]
2026
\end{center}

\newpage

% =====================================================================
%  Страницы 2+: содержательная часть задания
% =====================================================================
\thispagestyle{plain}
\singlespacing
\fontsize{12pt}{14pt}\selectfont
\noindent

% ===== Основная форма задания =====
\renewcommand{\arraystretch}{1.15}

% Заголовок раздела (по центру, жирно, на ширину всех трёх колонок).
\providecommand{\sectionheader}[1]{%
  \multicolumn{3}{|c|}{\rule{0pt}{5ex}\textbf{#1}\rule[-4ex]{0pt}{0pt}}}
% Тело раздела (на ширину всех трёх колонок).
\providecommand{\sectionbody}[1]{%
  \multicolumn{3}{|p{\dimexpr 1.8cm+3.5cm+10.7cm+4\tabcolsep+2\arrayrulewidth\relax}|}{#1}}

\begin{longtable}{|>{\centering\arraybackslash}p{1.8cm}|>{\centering\arraybackslash}p{3.5cm}|p{10.7cm}|}
\hline
\sectionheader{Цель работы} \\
\hline
\sectionbody{Разработка методики нейросетевого обнаружения
несанкционированных действий и аномальной активности в корпоративной
сети, валидируемой конечно-автоматным агентом активного тестирования с
библиотекой параметризованных шаблонов атак и блоком мониторинга
дрейфа распределений на основе индексов PSI, KL и MMD.} \\
\hline
\sectionheader{Решаемые задачи} \\
\hline
\sectionbody{%
1. Аналитический обзор подходов к обнаружению аномалий и
несанкционированных действий в сетевом трафике, сравнительный анализ
нейросетевых архитектур и публично доступных датасетов NIDS,
формализация функциональных и нефункциональных требований к системе.
\newline
2. Проектирование методики обнаружения: схема входного представления
flow-данных в скользящих окнах, обоснование выбора целевой
архитектуры детектора, схема калибровки порога под целевую частоту
ложных срабатываний и формирования алертов с уровнями критичности.
\newline
3. Разработка конечно-автоматного агента активного тестирования с
библиотекой параметризованных шаблонов сетевого трафика и инжектором
контролируемого дрейфа распределений.
\newline
4. Программная реализация прототипа на платформе Python\,/\,PyTorch с
интерфейсами FastAPI и Streamlit, юнит-тестирование критических
инвариантов pipeline и фиксация артефактов воспроизводимости.
\newline
5. Экспериментальная проверка эффективности методики: сравнение
целевой модели с базовыми подходами в идентичных условиях,
стресс-тестирование детектора в связке с агентом, оценка устойчивости
к контролируемому дрейфу распределений.
\newline
6. Оценка эксплуатационных характеристик прототипа и оформление
итогового отчёта, демонстрационных и презентационных материалов.} \\
\hline
\sectionheader{Объект исследования} \\
\hline
\sectionbody{Методы и технологии обнаружения аномальной активности и
несанкционированных действий в корпоративных сетях на основе
телеметрии сетевого трафика, включая статистические, машинно-обучающиеся
и нейросетевые подходы к классификации потоковых признаков.} \\
\hline
\sectionheader{Предмет исследования} \\
\hline
\sectionbody{Применение нейросетевых методов для выявления аномалий и
признаков несанкционированной активности по данным сетевого трафика, а
также разработка и применение программного агента на основе конечного
автомата с библиотекой параметризованных атакующих шаблонов для синтеза
трафика и сценариев активного тестирования системы обнаружения.} \\
\hline
\sectionheader{Актуальность работы} \\
\hline
\sectionbody{Актуальность темы работы обусловлена кратным ростом числа и
сложности компьютерных атак на отечественную информационную
инфраструктуру. По данным RED Security SOC, за первое полугодие
2025~года выявлено более 63 тысяч атак на российские организации, что
на 27\,\% превышает показатель аналогичного периода 2024~года.
Действующая нормативная база~--- Федеральный закон №\,187-ФЗ, приказы
ФСТЭК России №\,239 и №\,17, требования ГосСОПКА~--- предписывает
внедрение средств обнаружения вторжений на объектах КИИ. Параллельно
фиксируется дефицит специалистов SOC и потребность в импортонезависимых
программных средствах NIDS, что формирует прямой запрос на
методически прозрачные воспроизводимые программные прототипы с
поддержкой активного тестирования и мониторинга дрейфа распределений.} \\
\hline
\sectionheader{Теоретическая значимость} \\
\hline
\sectionbody{Сформирована единая воспроизводимая методика обнаружения
несанкционированных действий, объединяющая нейросетевой детектор на
основе гибридной архитектуры CNN-LSTM, конечно-автоматный агент
активного тестирования с библиотекой из семи параметризованных шаблонов
атак и одиннадцатью состояниями планирования режимов норма/атака/дрейф,
а также блок мониторинга дрейфа распределений на основе статистических
индексов PSI, KL и MMD. Методика формализована в виде шести этапов
протокола экспериментального исследования (E1--E6) с
bootstrap-доверительными интервалами 95\,\% и парным критерием
Уилкоксона; принципы переносимы на пары <<датасет + целевая модель>>,
отличные от UNSW-NB15/CNN-LSTM.} \\
\hline
\sectionheader{Практическая значимость} \\
\hline
\sectionbody{Реализован программный прототип на стеке
Python\,/\,PyTorch с интерфейсами FastAPI и Streamlit (порядка
4000~строк исходного кода, 11~модулей юнит-тестов, 86 проходящих тестов),
допускающий развёртывание в составе подсистемы NTA-аналитики
корпоративного SOC. На датасете UNSW-NB15 целевая архитектура CNN-LSTM
достигает $F_1\,=\,0{,}9735$ и $\mathrm{PR\text{-}AUC}\,=\,0{,}9945$
на F1-max пороге и $F_1\,=\,0{,}971$ при $\mathrm{FPR}\,=\,0{,}0095$ на
операционном пороге, что обеспечивает в одиннадцать раз меньшее число
ложных срабатываний по сравнению с XGBoost. Сформированы инженерные
рекомендации по интеграции с источниками NetFlow\,/\,IPFIX и калибровке
порогов под целевой эксплуатационный FPR.} \\
\hline
\sectionheader{Рекомендуемая литература} \\
\hline
\sectionbody{%
1. Goodfellow~I. Deep Learning / I.~Goodfellow, Y.~Bengio,
A.~Courville. --- Cambridge, MA: MIT Press, 2016. --- 800~p.
\newline
2. Khraisat~A. Survey of intrusion detection systems: techniques,
datasets and challenges / A.~Khraisat, I.~Gondal, P.~Vamplew,
J.~Kamruzzaman // Cybersecurity. --- 2019. --- Vol.~2, No.~20. ---
DOI 10.1186/s42400-019-0038-7.
\newline
3. Ferrag~M.A. Deep learning for cyber security intrusion detection:
Approaches, datasets, and comparative study / M.A.~Ferrag,
L.~Maglaras, S.~Moschoyiannis, H.~Janicke // Journal of Information
Security and Applications. --- 2020. --- Vol.~50. ---
DOI 10.1016/j.jisa.2019.102419.
\newline
4. Moustafa~N. UNSW-NB15: a comprehensive data set for network
intrusion detection systems / N.~Moustafa, J.~Slay //
Military Communications and Information Systems Conference (MilCIS),
Canberra, Australia. --- 2015. --- P.~1--6. ---
DOI 10.1109/MilCIS.2015.7348942.} \\
\hline
\sectionheader{Календарный план} \\
\hline
\textbf{Номер этапа} & \textbf{Срок выполнения} & \multicolumn{1}{c|}{\textbf{Наименование этапа}} \\
\hline
1 & 11.05.2026 -- 15.05.2026 &
Аналитико-постановочный этап: аналитический обзор предметной области,
сравнительный анализ нейросетевых архитектур и датасетов NIDS,
формализация функциональных и нефункциональных требований \\
\hline
2 & 16.05.2026 -- 20.05.2026 &
Проектный этап: проектирование методики обнаружения, архитектуры
детектора, конечно-автоматного агента и протокола экспериментального
исследования \\
\hline
3 & 21.05.2026 -- 30.05.2026 &
Реализационный этап: программная реализация прототипа системы,
конечно-автоматного агента и подсистемы мониторинга дрейфа
распределений \\
\hline
4 & 31.05.2026 -- 04.06.2026 &
Экспериментально-аналитический этап: проверка эффективности методики,
стресс-тестирование, оценка устойчивости к контролируемому дрейфу
распределений \\
\hline
5 & 05.06.2026 -- 07.06.2026 &
Заключительный этап: оценка эксплуатационных характеристик,
оформление итогового отчёта, демонстрационных и презентационных
материалов \\
\hline
\end{longtable}

\vspace{12pt}

\noindent
\textbf{Дата выдачи задания:} 11.05.2026 \\
\textbf{Срок выполнения работы:} 07.06.2026

\vspace{30pt}

\noindent
\begin{tabular}{@{}p{0.45\textwidth}@{\hspace{0.5cm}}p{0.45\textwidth}@{}}
Руководитель ВКР & Студент \\[28pt]
\makebox[3.5cm]{\hrulefill}\,/\,Когос К.Г./ & \makebox[3.5cm]{\hrulefill}\,/\,Титов Д.И./ \\
\centering(подпись, ФИО) & \centering(подпись, ФИО) \tabularnewline
\end{tabular}

\onehalfspacing
